\section{Approach: Fusing Log-to-Code Retrieval with Incident Memory}
\label{sec:approach}

% OUTLINE (FIGURE 1: architecture -- new drawing, /paper-diagrams style):
%
% 3.1 Overview. Input per the challenge: issue text + logs + snapshot.
%   Output: ranked files, affected module, (explanation -- see status in
%   Sec. discussion). Five ranked signals over the snapshot's source files,
%   fused by weighted reciprocal-rank fusion.
%
% 3.2 Signals (map each to the challenge's three key difficulties):
%   S1 trace mentions: explicit file mentions (Foo.java, foo.cc:123, JVM
%      frames) + stem fallback. [also serves as the mandatory baseline]
%   S2 path retrieval: BM25 over path tokens.
%   S3 content retrieval: BM25 over file content of a candidate pool;
%      document frequencies computed inside the pool (no cross-snapshot
%      leakage). [evidence-to-code semantic retrieval]
%   S4+S5 incident memory: past incidents ranked by text similarity
%      (depth 20) and by the incident retriever; each retrieved incident
%      votes 1/rank for the files its fix touched, mapped into the current
%      snapshot. [root cause beyond symptoms: recurring failure modes
%      concentrate in the same code]
%
% 3.3 The memory firewall (what keeps the signal honest): exclude the
%   query issue, its duplicate cluster, any incident resolved after the
%   query's creation time, any incident sharing the query's fix commit.
%   Every exclusion counted and reported.
%
% 3.4 Regime-conditioned fusion. The easy/hard stratum is gold-derived and
%   unusable at inference; "does the text contain an explicit file
%   mention" IS observable. Weights tuned per regime on a time-based dev
%   split (earliest 30\%), applied unchanged to the later 70\%.
